\documentclass[11pt]{article}
\usepackage[utf8]{inputenc}
\usepackage[T1]{fontenc}
\usepackage{amsmath}
\usepackage{amssymb}
\usepackage{multicol}
\usepackage{geometry}
\usepackage{enumitem}
\usepackage{xcolor}
\usepackage{titlesec}

% Layout
\geometry{a4paper, left=1cm, right=1cm, top=0.8cm, bottom=1.2cm}
\setlength{\columnseprule}{0.4pt}

% Títulos
\titleformat{\section}{\normalfont\Large\bfseries\color{blue}}{}{0em}{}
\titleformat{\subsection}{\normalfont\large\bfseries\color{red}}{}{0em}{}

\title{\textcolor{red}{Segundo Trimestre: Matemática\\Equação do 1º Grau – Questões Aplicadas}}
\author{Professor: Jefferson}
\date{}

\begin{document}

\maketitle
\vspace{-1cm}

\begin{multicols}{2}

\section*{Questões Aplicadas}

\subsection*{Questão 1}
Uma loja cobra R\$ 20 de taxa fixa mais R\$ 15 por produto. Se o total pago foi R\$ 95, quantos produtos foram comprados?

\textbf{Comentário:}
\[
20 + 15x = 95
\Rightarrow 15x = 75
\Rightarrow x = 5
\]

\textbf{Resposta:} 5 produtos

\subsection*{Questão 2}
O triplo da idade de Pedro menos 9 é igual a 30. Qual é a idade de Pedro?

\textbf{Comentário:}
\[
3x - 9 = 30
\Rightarrow 3x = 39
\Rightarrow x = 13
\]

\textbf{Resposta:} 13 anos

\subsection*{Questão 3}
Uma corrida de táxi cobra R\$ 6 de bandeirada mais R\$ 4 por quilômetro. Se o valor total foi R\$ 46, quantos quilômetros foram percorridos?

\textbf{Comentário:}
\[
6 + 4x = 46
\Rightarrow 4x = 40
\Rightarrow x = 10
\]

\textbf{Resposta:} 10 km

\subsection*{Questão 4}
A soma de um número com sua metade é igual a 27. Qual é esse número?

\textbf{Comentário:}
\[
x + \frac{x}{2} = 27
\Rightarrow \frac{3x}{2} = 27
\Rightarrow x = 18
\]

\textbf{Resposta:} 18

\subsection*{Questão 5}
Um estacionamento cobra R\$ 12 mais R\$ 6 por hora. Se o total pago foi R\$ 54, quantas horas o carro ficou estacionado?

\textbf{Comentário:}
\[
12 + 6x = 54
\Rightarrow 6x = 42
\Rightarrow x = 7
\]

\textbf{Resposta:} 7 horas

\subsection*{Questão 6}
O dobro de um número somado com 14 resulta em 50. Qual é esse número?

\textbf{Comentário:}
\[
2x + 14 = 50
\Rightarrow 2x = 36
\Rightarrow x = 18
\]

\textbf{Resposta:} 18

\subsection*{Questão 7}
A diferença entre um número e 11 é igual a 29. Qual é esse número?

\textbf{Comentário:}
\[
x - 11 = 29
\Rightarrow x = 40
\]

\textbf{Resposta:} 40

\subsection*{Questão 8}
O perímetro de um retângulo é 70 cm. Sabendo que a largura mede 15 cm, qual é o comprimento?

\textbf{Comentário:}
\[
2(x + 15) = 70
\Rightarrow x + 15 = 35
\Rightarrow x = 20
\]

\textbf{Resposta:} 20 cm

\subsection*{Questão 9}
Uma pessoa recebe R\$ 200 fixos mais R\$ 40 por dia trabalhado. Se recebeu R\$ 520, quantos dias trabalhou?

\textbf{Comentário:}
\[
200 + 40x = 520
\Rightarrow 40x = 320
\Rightarrow x = 8
\]

\textbf{Resposta:} 8 dias

\subsection*{Questão 10}
Um número somado com seu dobro é igual a 63. Qual é esse número?

\textbf{Comentário:}
\[
x + 2x = 63
\Rightarrow 3x = 63
\Rightarrow x = 21
\]

\textbf{Resposta:} 21

\subsection*{Questão 11}
Uma escola arrecadou R\$ 1.500 vendendo ingressos a R\$ 25 cada. Quantos ingressos foram vendidos?

\textbf{Comentário:}
\[
25x = 1500
\Rightarrow x = 60
\]

\textbf{Resposta:} 60 ingressos

\subsection*{Questão 12}
O quádruplo de um número menos 16 é igual a 48. Qual é esse número?

\textbf{Comentário:}
\[
4x - 16 = 48
\Rightarrow 4x = 64
\Rightarrow x = 16
\]

\textbf{Resposta:} 16

\subsection*{Questão 13}
Um plano de celular custa R\$ 60 mais R\$ 5 por GB utilizado. Se a conta foi R\$ 85, quantos GB foram usados?

\textbf{Comentário:}
\[
60 + 5x = 85
\Rightarrow 5x = 25
\Rightarrow x = 5
\]

\textbf{Resposta:} 5 GB

\subsection*{Questão 14}
A metade de um número somada com 9 é igual a 33. Qual é esse número?

\textbf{Comentário:}
\[
\frac{x}{2} + 9 = 33
\Rightarrow \frac{x}{2} = 24
\Rightarrow x = 48
\]

\textbf{Resposta:} 48

\subsection*{Questão 15}
Após um desconto de R\$ 35, um produto passou a custar R\$ 95. Qual era o preço original?

\textbf{Comentário:}
\[
x - 35 = 95
\Rightarrow x = 130
\]

\textbf{Resposta:} R\$ 130

\subsection*{Questão 16}
O triplo de um número é igual a esse número somado com 36. Qual é o número?

\textbf{Comentário:}
\[
3x = x + 36
\Rightarrow 2x = 36
\Rightarrow x = 18
\]

\textbf{Resposta:} 18

\subsection*{Questão 17}
Uma sala tem o dobro de alunos de outra. Sabendo que a menor sala tem 22 alunos, quantos alunos há na maior?

\textbf{Comentário:}
\[
x = 2 \cdot 22
\]

\textbf{Resposta:} 44 alunos

\subsection*{Questão 18}
Um número dividido por 6 resulta em 14. Qual é esse número?

\textbf{Comentário:}
\[
\frac{x}{6} = 14
\Rightarrow x = 84
\]

\textbf{Resposta:} 84

\subsection*{Questão 19}
Uma mensalidade custa R\$ 100 mais R\$ 20 por aula extra. Se o total pago foi R\$ 180, quantas aulas extras foram feitas?

\textbf{Comentário:}
\[
100 + 20x = 180
\Rightarrow 20x = 80
\Rightarrow x = 4
\]

\textbf{Resposta:} 4 aulas

\subsection*{Questão 20}
A soma de dois números consecutivos é 89. Qual é o menor deles?

\textbf{Comentário:}
\[
x + (x + 1) = 89
\Rightarrow 2x = 88
\Rightarrow x = 44
\]

\textbf{Resposta:} 44

\end{multicols}

\end{document}

